\section{Basic Syntax}

\subsection{Basic syntax}

The runnable version of this chapter lives in
\href{/BasicSyntax-17b3beb7abc501549129aaa1412511b4.lean}{\texttt{LearningLean/BasicSyntax.lean}}.

\subsubsection{Comments}

\begin{verbatim}
- - This is a single -line comment.

/ -
This is a multi -line comment.
Here is another line.
 -/
\end{verbatim}

\subsubsection{Inspecting terms and types with \texttt{\#check}}

\texttt{\#check} asks Lean for the type of an expression without evaluating it.

\begin{verbatim}
#check true          - - Bool
#check 42            - - Nat
#check "Hello, World!" - - String
#check Type          - - Type 1
#check Nat           - - Type
#check ['h', 'e', 'l', 'l', 'o'] - - List Char
\end{verbatim}

\subsubsection{Printing definitions with \texttt{\#print}}

\texttt{\#print} shows how something is defined.

\begin{verbatim}
#print Bool
\end{verbatim}

\texttt{Type} itself cannot be printed with \texttt{\#print}, because it is a type rather than a
definition.

\subsubsection{Definitions and functions}

\begin{verbatim}
def pi : Float := 3.1415926

#check pi

 - - An anonymous function (lambda) taking two natural numbers.
#check \lambda (a b : Nat) => a + b
\end{verbatim}